\section{Changing of the Guard}

The third dimension of history understood as the impact of imaginal world on human history. Wars, migrations, expansions, conversions, and so on, need to be seen in the light of this higher dimension, above and beyond the material, biological, economic, and ideological causes.

We have seen that \textbf{Julius Evola} attributes the movements of human events to hidden, transcendental forces. However, he gives no details about those forces. If we turn to \textbf{Rene Guenon}, we find an acknowledgment of those forces, but treated in a superficial way. Guenon distinguishes between the Unmanifested, Formless Manifestation, and Formal Manifestation.

Formal Manifestation is our human state, body and soul, along with the existing world. Guenon's writings are oriented to pure metaphysics, transcending all states, what he calls the \textbf{Supreme Identity} from Sufi teachings. Formless manifestation includes beings that are not individuated in the material world. As part of manifestation, its study is not properly metaphysics, but involves the traditional sciences, e.g., cosmology. This is not Guenon's task as he sees it.

There are two ways to look at formless manifestation. From the strictly human perspective, it is transcendent to the human state. However, from the perspective of the Being, the human state is just one possible state, among many other higher states of formless manifestation. Thus, in the Multiple States of the Being, Guenon mentions the angels, which are seen as exterior beings in exoteric religion. However, esoterically, they are higher states. These states are between the Absolute and the human state. For details of those states, we need to look elsewhere.

Please don't think of ``angels" in the merely sentimental way that you may be accustomed to. We mean, by ``angels", those superior Intelligences which form part of formless manifestation. Just focusing on one part of hidden history, namely, the history of nations and races, we find this from \textbf{Thomas Aquinas}:

\begin{quotex}
For in human affairs there is a common good which is, in fact, the good of a state or a people, and this seems to belong to the order of Principalities. … the arrangement of kingdoms and the changing of domination from one people to another ought to belong to the ministry of this order. Also, the instruction of those who occupy the position of leaders among men concerning matters pertinent to the administration of their rule seems to be the concern of this order. \flright{\textsc{Thomas Aquinas}, \emph{Summa contra Gentiles}, III, 80}

\end{quotex}
So there is an order of the angels, called Principalities, who influence or direct the various nations, ethnicities, and races. Wars, migrations, expansions, conversions, and so on, need to be seen in the light of this higher dimension, above and beyond the material, biological, economic, and ideological causes.

Book III of the Summa deals only with topics that can be known solely from human experience and reason, and do not depend on any special revelation. So what evidence is there to indicate that human events are subject to hidden, transcendental forces? We can propose these three:

\begin{enumerate}
\item The rapidity with which certain ideas and intellectual movements can grab hold of a nation, or particular segments of a nation. Seldom are they the result of intelligent forethought, but instead seem to arise spontaneously. 
\item The actual results of revolutions or political programs do not comport with the original intentions. Even if this is not noticed right away, then certainly over time, any initial successes will be followed by stages of decline, until the final result is the complete opposite of the original movement. 
\item It was the common experience of traditional cultures. 
\end{enumerate}
The first two are clear and obvious enough, so we can focus on the third. We recently posted an article by \textbf{Guido de Giorgio} about the recovery of the past. As he writes:

\begin{quotex}
Whoever intends to remain in the pure domain of traditional truth, always turns, logically, toward the past, to retrace the stages of certitude and add them to his experience. 

\end{quotex}
So, if we can no longer experience the world as our ancestors once did, we need to retrace our steps. The way to start is a form of Hermetic Meditation, what \textbf{Henry Corbin} calls creative imagination. In the Middle Ages, the imagination was counted as one of the inner wits.

As an example, in the Trojan War, the gods and goddesses took different sides in the conflict. In the \emph{Convivio}, \textbf{Dante} relates this to the higher Intelligences, or principalities. He writes:

\begin{quotex}
There are others like the eminent Plato who maintain that not only are there as many Intelligences as there are spheres in heaven, but also as many as there are species of things, for example one for men, another for gold, another for dimension, and so on. He held that just as the heavenly Intelligences each brought their sphere into being, so other Intelligences brought into being all other things and exemplars, each in its own species; and Plato called them Ideas, that is to say universal forms and natures. … The pagans called them God and Goddesses. 

\end{quotex}
Dante is saying much more here than that these higher intelligences are involved in human affairs in terms of politics or war. The gods and goddesses also represented certain qualities: beauty, wisdom, martial spirit, and so on, what Plato called the Ideas. In Guenon's scheme, the ideas are possibilities in the Infinity of the unmanifest Absolute. However, they have another existence in the imaginal world of formless manifestation, where the ideas are experienced not as abstractions, but as living beings.

As De Giorgio pointed out, we cannot simply regard the past as simply of antiquarian interest, as something exterior to us. Rather, traditional truths must be ``found again and vivified". Specifically, one must live them as these higher states of being. As Evola never tires of pointing out, the way of action leads to the awareness and achievement of these higher states. Effective action, or the recovery of Tradition, can only arise out of these higher states.

\flrightit{Posted on 2012-06-14 by Cologero }
